\section{Introduction}

The hot-spots problem arose from Rauch's discussion of the long-time behavior of heat flow with insulating boundary conditions \cite{Rauch1975}. Its spectral formulation asks whether the extrema of a first-positive Neumann eigenfunction occur on the boundary. Bañuelos and Burdzy distinguished strict and nonstrict formulations, and assertions about every eigenfunction from the existence of one favorable eigenfunction \cite{BanuelosBurdzy1999}. These distinctions matter when the eigenvalue is multiple.

Let $\Omega\subset\R^2$ be a nonempty bounded simply connected open set with $C^\infty$ boundary. In particular, $\Omega$ is connected. Set
\begin{equation}\label{eq:mu}
\mu=\mu_1(\Omega)
=\inf_{\substack{0\ne v\in H^1(\Omega)\\ \int_\Omega v=0}}
\frac{\int_\Omega|\grad v|^2}{\int_\Omega|v|^2},
\end{equation}
and let $\V$ be the real Neumann eigenspace for $\mu$. We count only positive eigenvalues in this notation; sources that count the constant eigenfunction first write $\mu_2$ for this same eigenvalue.

\begin{theorem}\label{thm:main}
For every $0\ne u\in\V$,
\[
\grad u(x)\ne0\qquad(x\in\Omega).
\]
In particular,
\begin{equation}\label{eq:hotspots}
\min_{y\in\partial\Omega}u(y)<u(x)<\max_{y\in\partial\Omega}u(y)
\qquad(x\in\Omega).
\end{equation}
\end{theorem}

The theorem imposes no convexity or symmetry condition. It applies to every member of the first positive eigenspace, including when the eigenvalue is multiple. The nonvanishing-gradient conclusion is stronger than \eqref{eq:hotspots}, since it also rules out interior saddle points.

\paragraph{Background and methods.}
Topology matters in the hot-spots problem. Burdzy and Werner constructed a planar counterexample with two holes and a simple first-positive eigenvalue \cite{BurdzyWerner1999}. Bass and Burdzy obtained examples with both extrema strictly in the interior \cite{BassBurdzy2000}, and Burdzy subsequently obtained this behavior with just one hole \cite{Burdzy2005}. The simply connected planar conjecture is stated explicitly in \cite[Conjecture~1.2(ii)]{Burdzy2005}.

Important geometric special cases already establish strong gradient conclusions. Jerison and Nadirashvili proved a positive directional derivative for every first-positive eigenfunction on convex planar domains with two orthogonal reflection symmetries, allowing multiple eigenvalues \cite[Theorem~1.4]{JerisonNadirashvili2000}. Probabilistic coupling methods developed by Bañuelos and Burdzy \cite{BanuelosBurdzy1999} led to further results. Pascu proved boundary-only extrema for antisymmetric eigenfunctions on $C^{1,\alpha}$ convex planar domains with one reflection symmetry, where $0<\alpha<1$ \cite{Pascu2002}. Atar and Burdzy treated every first-positive eigenfunction on lip domains, a class of Lipschitz domains bounded between graphs of functions with Lipschitz constants at most one \cite{AtarBurdzy2004}. Miyamoto obtained interior critical-point exclusion under a spectral-diameter condition, giving nonsymmetric nearly circular convex examples \cite[Lemma~1.2 and Theorem~A]{Miyamoto2007}.

The polygonal problem has a complementary development. Siudeja proved hot-spots results for a class of acute triangles \cite{Siudeja2015}. Judge and Mondal, with their subsequent erratum, excluded interior critical points for every first-positive Neumann eigenfunction on every Euclidean triangle \cite{JudgeMondal2020,JudgeMondal2022}. Chen, Gui and Yao established the finer classification of nonvertex critical points \cite{ChenGuiYao2026}. These results concern domains with corners; \Cref{thm:main} concerns the full smooth simply connected planar class.

Recent quantitative work measures how far the hot-spots conclusion can fail. For general bounded connected Lipschitz domains, de Dios Pont, Hsu and Taylor determined the sharp dimension-dependent upper bound for the ratio of a first-positive Neumann eigenfunction's supremum over the domain to its maximum on the boundary \cite[Definition~1 and Theorem~4]{deDiosPontHsuTaylor2025}. Deng, Jiang and Yang obtained quantitative bounds for convex domains in two-dimensional space forms \cite{DengJiangYang2026}. Such ratios concern extrema; excluding every interior critical point also requires excluding saddle points.

Our starting point is Rohleder's variational principle for tangent vector fields \cite{Rohleder2021,Rohleder2024}. In its lip-domain application, taking componentwise absolute values after a suitable rotation preserves the required tangent boundary condition and yields directional monotonicity. We use scalar boundary multipliers to keep tangency on an arbitrary smooth boundary, and construct those multipliers through a reciprocal Green-kernel theorem.

The difficulty is to control the gradient without a preferred direction in the domain. Each Cartesian derivative of an eigenfunction solves the same Helmholtz equation, but its boundary values need not have a fixed sign. The Neumann condition instead constrains the two derivatives together: their vector is tangent to the boundary. We use that constraint through a nonnegative quadratic form whose nullspace consists exactly of first-positive eigenfunction gradients. The remaining task is to construct enough scalar boundary multipliers that preserve this nullspace if an interior critical point exists.

\paragraph{The mechanism.}
Write $\D=\{z\in\C:|z|<1\}$ and $\Sone=\partial\D$, and let $\Phi:\D\to\Omega$ be a smooth conformal parametrization. The proof begins with the variational principle for tangent vector fields associated with Neumann gradients. For a tangent vector field $X$, its divergence--curl energy is at least $\mu\norm{X}_2^2$, with equality exactly for gradients of functions in $\V$. After this conformal change of variables, solving the scalar Helmholtz equation in each Cartesian component gives a nonnegative quadratic form $E$ on boundary vector data. If $g$ is the boundary gradient of an eigenfunction, then $E(g)=0$, and a scalar multiplier $b$ satisfies
\[
E(bg)=\frac12\iint_{\Sone\times\Sone}
N(s,t)(b(s)-b(t))^2g(s)\cdot g(t)\dd s\dd t.
\]
Here $N$ is a positive boundary kernel, and the corresponding interior evaluation kernel is denoted by $K$.

The sign of $g(s)\cdot g(t)$ prevents a direct positivity argument in this identity. The main kernel result supplies a family of multipliers whose squared differences remove the factor $N$. For every interior point $p$,
\[
D_p(s,t)=\frac{K(p,s)K(p,t)}{N(s,t)}\quad(s\ne t),
\qquad D_p(s,s)=0,
\]
is conditionally negative semidefinite. It is therefore the squared distance of an injective Lipschitz map from the circle to a real Hilbert space. Applying the preceding energy identity to its coordinates gives
\[
\sum_j E(b_jg)=\frac12|\grad u(\Phi(p))|^2.
\]
A critical point would thus produce too many eigenfunction gradients, contradicting the planar multiplicity bound or boundary uniqueness.

The construction of $D_p$ is the central analytic step. We first regularize the disk Green function so that the entrywise reciprocal of every finite kernel matrix has at most one positive eigenvalue. This property survives successive rank-one updates along columns of the current kernel matrix, which sum the Green resolvent over a finite set of integration nodes. Positive quadrature and monotone limits then pass from those matrices to the singular Green kernel and finally to its boundary kernels. Taking the boundary limit before removing a compact-support cutoff avoids treating a Poisson kernel as an $L^2$ function when it need not be one.

The kernel theorem, \Cref{thm:negative-kernel}, holds for every bounded nonnegative disk potential satisfying the subcritical Dirichlet inequality \eqref{eq:subcritical}, independently of the conformal origin of that potential. The reciprocal-matrix condition comes from the McCullough--Quiggin characterization of complete Nevanlinna--Pick kernels \cite{McCullough1992,McCullough1994,Quiggin1993,Quiggin1994}; see also \cite[Corollary~1.12]{AglerMcCarthy2000}. Schoenberg's squared-distance correspondence supplies the Hilbert-space realization \cite{Schoenberg1938}. The analytic work here concerns the singular Green kernel, its positive-potential resolvent, and its boundary limit; the resulting theorem can be used independently of the eigenfunction argument.

\paragraph{Organization and background.}
\Cref{sec:preliminaries} proves the spectral gap, the smooth conformal reduction, and the multiplicity bound. \Cref{sec:vector} proves the vector variational principle. \Cref{sec:boundary} constructs the boundary kernels and derives the multiplier identity. \Cref{sec:reciprocal} proves their conditional negativity, and \Cref{sec:conclusion} completes the argument.
We use standard Sobolev compactness, Poincar\'e and trace inequalities on smooth bounded domains, Poisson solvability and elliptic boundary regularity, the maximum and boundary point principles, elementary complex analysis, and planar separation. All specialized ingredients used in the proof are established below. The smooth-boundary regularity input is \cite[Chapter~III, Theorems~6.1, 6.4 and~6.5]{Showalter1994}: $L^2$ Poisson data give $H^2$ solutions under homogeneous Dirichlet or Neumann conditions, and higher regularity follows with smoother data. Subtracting a smooth extension handles the inhomogeneous Dirichlet data used below. Bootstrap and Sobolev embedding make weak Neumann eigenfunctions smooth up to the boundary. Interior solutions of $(\Delta+\mu)v=0$ are analytic; alternatively, $e^{\sqrt\mu t}v(x)$ is harmonic in one additional variable.
